\documentclass{article}

\usepackage[utf8]{inputenc}
\usepackage[T1]{fontenc}
\usepackage[a4paper,margin=2.5cm]{geometry}
\usepackage{xcolor}
\usepackage{amsmath}
\usepackage{amssymb}
\usepackage{listings}

\lstset{
  basicstyle=\ttfamily\small,
  breaklines=true,
  frame=single,
  numbers=left,
  numberstyle=\tiny\color{gray},
  keywordstyle=\bfseries\color{blue!60!black},
  commentstyle=\itshape\color{green!40!black},
  stringstyle=\color{red!50!black},
  showstringspaces=false,
  tabsize=2
}

\title{Astrazione dei Dati e Paradigma Orientato agli Oggetti}
\author{}
\date{}

\begin{document}

\maketitle

\section{Astrazione dei Dati}

La macchina fisica manipola dati di un solo tipo fondamentale: stringhe di bit. I \textbf{tipi} di un linguaggio di alto livello organizzano questo universo indifferenziato associando ai valori una forma di incapsulamento. Un tipo può essere visto come una \textbf{capsula} che:
\begin{itemize}
  \item raccoglie un insieme di valori;
  \item stabilisce quali operazioni possono manipolare tali valori;
  \item nasconde, nei linguaggi \textit{type-safe}, la rappresentazione concreta.
\end{itemize}

Un linguaggio che offre solo array, record, puntatori e tipi ricorsivi permette di costruire strutture composte, ma non sempre consente al programmatore di definire nuovi tipi allo stesso livello di astrazione dei tipi primitivi. Il problema centrale è quindi fornire meccanismi di \textbf{data abstraction}, cioè strumenti per separare l'uso di una struttura dalla sua rappresentazione interna.

\subsection{Limite dei Tipi Compositi}

Una possibile implementazione di una pila di interi in uno pseudo-linguaggio è la seguente:

\begin{lstlisting}
type Int_Stack = struct{
    int P[100];    // the stack proper
    int top;       // first readable element
}

Int_Stack create_stack(){
    Int_Stack s = new Int_Stack();
    s.top = 0;
    return s;
}

Int_Stack push(Int_Stack s, int k){
    if (s.top == 100) error;
    s.P[s.top] = k;
    s.top = s.top + 1;
    return s;
}

int top(Int_Stack s){
    return s.P[s.top];
}

Int_Stack pop(Int_Stack s){
    if (s.top == 0) error;
    s.top = s.top - 1;
    return s;
}

bool empty(Int_Stack s){
    return (s.top == 0);
}
\end{lstlisting}

Questa definizione suggerisce che una pila debba essere manipolata solo tramite \texttt{create\_stack}, \texttt{push}, \texttt{top}, \texttt{pop} ed \texttt{empty}. Tuttavia la rappresentazione è esposta: un client può accedere direttamente all'array interno.

\begin{lstlisting}
int second_from_top(Int_Stack s){
    return s.P[s.top - 1];
}
\end{lstlisting}

L'accesso diretto viola l'astrazione perché il client dipende dalla rappresentazione concreta della pila.

\section{Abstract Data Types}

Un \textbf{abstract data type} (\textbf{ADT}) è un meccanismo linguistico che consente di definire un nuovo tipo rendendo opaca la sua rappresentazione.

\subsection{Componenti di un ADT}

Un ADT comprende:
\begin{enumerate}
  \item un nome per il tipo astratto;
  \item una rappresentazione concreta del tipo;
  \item una \textbf{signature} o \textbf{interface}, cioè i nomi e i tipi delle operazioni esportate;
  \item l'implementazione delle operazioni;
  \item una capsula di sicurezza che separa l'interfaccia dall'implementazione.
\end{enumerate}

Formalmente, un ADT può essere rappresentato come:
\[
\mathcal{A} = (T_{\mathrm{abs}}, T_{\mathrm{rep}}, \Sigma, I)
\]
dove $T_{\mathrm{abs}}$ è il nome astratto, $T_{\mathrm{rep}}$ è il tipo di rappresentazione, $\Sigma$ è la signature pubblica e $I$ è l'implementazione privata.

\subsection{Esempio: ADT per Stack di Interi}

\begin{lstlisting}
abstype Int_Stack{
    type Int_Stack = struct{
        int P[100];
        int n;
        int top;
    }

    signature
        Int_Stack create_stack();
        Int_Stack push(Int_Stack s, int k);
        int top(Int_Stack s);
        Int_Stack pop(Int_Stack s);
        bool empty(Int_Stack s);

    operations
        Int_Stack create_stack(){
            Int_Stack s = new Int_Stack();
            s.n = 0;
            s.top = 0;
            return s;
        }

        Int_Stack push(Int_Stack s, int k){
            if (s.n == 100) error;
            s.n = s.n + 1;
            s.P[s.top] = k;
            s.top = s.top + 1;
            return s;
        }

        int top(Int_Stack s){
            return s.P[s.top];
        }

        Int_Stack pop(Int_Stack s){
            if (s.n == 0) error;
            s.n = s.n - 1;
            s.top = s.top - 1;
            return s;
        }

        bool empty(Int_Stack s){
            return (s.n == 0);
        }
}
\end{lstlisting}

All'interno della definizione, \texttt{Int\_Stack} è sinonimo della sua rappresentazione concreta; all'esterno, invece, non esiste alcuna relazione accessibile tra il tipo astratto e il record che lo implementa. Un client non può applicare selettori di campo a un valore di tipo \texttt{Int\_Stack}.

\subsection{Interfaccia e Implementazione}

Un ADT è una capsula opaca:
\begin{itemize}
  \item all'esterno sono visibili il nome del tipo e le operazioni esportate;
  \item all'interno sono presenti la rappresentazione e il codice delle operazioni;
  \item ogni accesso alla rappresentazione deve passare attraverso l'interfaccia.
\end{itemize}

La \textbf{signature} è quindi il contratto sintattico dell'ADT, mentre l'implementazione è nascosta.

\section{Information Hiding}

Il principio di \textbf{information hiding} separa l'uso di un componente dalla sua definizione interna. Per le funzioni, tale separazione riguarda il corpo della funzione; per gli ADT riguarda sia il codice delle operazioni sia la rappresentazione dei dati.

\subsection{Costruttori, Trasformatori e Osservatori}

Le operazioni di un ADT $T$ si dividono concettualmente in:
\begin{itemize}
  \item \textbf{costruttori}: costruiscono nuovi valori di tipo $T$;
  \item \textbf{trasformatori}: producono valori di tipo $T$ a partire da valori già esistenti;
  \item \textbf{osservatori}: producono valori di tipi diversi da $T$ osservando valori di tipo $T$.
\end{itemize}

Se un trasformatore ha tipo:
\[
t : S_1 \times \cdots \times S_k \to T
\]
allora, per ogni tupla di argomenti $(s_1,\ldots,s_k)$, il valore $t(s_1,\ldots,s_k)$ deve essere costruibile tramite i costruttori dell'ADT.

Nel caso della pila:
\begin{itemize}
  \item \texttt{create\_stack} e \texttt{push} sono costruttori;
  \item \texttt{pop} è un trasformatore;
  \item \texttt{top} ed \texttt{empty} sono osservatori.
\end{itemize}

\subsection{ADT Implementato tramite un Altro ADT}

\begin{lstlisting}
abstype Int_Var{
    type Int_Var = Int_Stack;

    signature
        Int_Var create_var();
        int deref(Int_Var v);
        Int_Var assign(Int_Var v, int n);

    operations
        Int_Var create_var(){
            return push(create_stack(), 0);
        }

        int deref(Int_Var v){
            return top(v);
        }

        Int_Var assign(Int_Var v, int n){
            return push(pop(v), n);
        }
}
\end{lstlisting}

L'implementazione di \texttt{Int\_Var} usa l'interfaccia di \texttt{Int\_Stack}, non la sua rappresentazione interna. Questo mostra che gli ADT possono essere composti mantenendo l'opacità.

\section{Indipendenza dalla Rappresentazione}

La \textbf{representation independence} afferma che due implementazioni corrette della stessa specifica di un ADT sono indistinguibili dai client che usano solo l'interfaccia.

\subsection{Seconda Implementazione dello Stack}

\begin{lstlisting}
abstype Int_Stack{
    type Int_Stack = struct{
        int info;
        Int_Stack next;
    }

    signature
        Int_Stack create_stack();
        Int_Stack push(Int_Stack s, int k);
        int top(Int_Stack s);
        Int_Stack pop(Int_Stack s);
        bool empty(Int_Stack s);

    operations
        Int_Stack create_stack(){
            return null;
        }

        Int_Stack push(Int_Stack s, int k){
            Int_Stack tmp = new Int_Stack();
            tmp.info = k;
            tmp.next = s;
            return tmp;
        }

        int top(Int_Stack s){
            return s.info;
        }

        Int_Stack pop(Int_Stack s){
            return s.next;
        }

        bool empty(Int_Stack s){
            return (s == null);
        }
}
\end{lstlisting}

La rappresentazione è ora una lista concatenata, non più un array. Un client corretto, che usa soltanto la specifica astratta, non osserva differenze.

\subsection{Specificazione Astratta}

Una possibile specifica dell'ADT \texttt{Int\_Stack} è:
\begin{itemize}
  \item \texttt{create\_stack} crea una pila vuota;
  \item \texttt{push} inserisce un elemento;
  \item \texttt{top} restituisce l'elemento in cima senza modificare la pila;
  \item \texttt{pop} rimuove l'elemento in cima;
  \item \texttt{empty} è vero se e solo se la pila è vuota.
\end{itemize}

La forma forte del principio è:
\[
I_1 \models \mathrm{Spec}
\land
I_2 \models \mathrm{Spec}
\Rightarrow
I_1 \equiv_{\mathrm{obs}} I_2
\]
dove $I_1$ e $I_2$ sono implementazioni corrette della stessa specifica e $\equiv_{\mathrm{obs}}$ indica indistinguibilità osservazionale per i client.

La forma debole, garantita dal type checking in linguaggi \textit{type-safe} con ADT, afferma che sostituire un'implementazione con un'altra avente la stessa signature non introduce errori di tipo.

\section{Moduli}

Gli ADT sono orientati alla \textbf{programmazione in piccolo}: incapsulano un tipo e le sue operazioni. I \textbf{moduli} o \textbf{package} sono invece meccanismi per la \textbf{programmazione in grande}: raccolgono più tipi, variabili e funzioni, controllando la visibilità delle dichiarazioni.

\subsection{Caratteristiche Linguistiche}

Un modulo:
\begin{itemize}
  \item ha una parte \textbf{pubblica}, visibile ai client;
  \item ha una parte \textbf{privata}, invisibile all'esterno;
  \item può importare altri moduli;
  \item può esportare tipi astratti, variabili e operazioni;
  \item può essere generico, cioè parametrico rispetto a tipi.
\end{itemize}

\subsection{Esempio di Moduli}

\begin{lstlisting}
module Buffer imports Counter{
    public
        type Buf;
        void insert(reference Buf f, int n);
        int get(Buf b);
        Count c; // how many times buffer has been used

    private imports Queue{
        type Buf = Queue;

        void insert(reference Buf b, int n){
            inqueue(b, n);
            inc(c);
        }

        int get(Buf b){
            return dequeue(b);
            inc(c);
        }

        init_counter(c); // module initialisation part
    }
}

module Counter{
    public
        type Count;
        void init_counter(reference Count c);
        int get(Count c);
        void inc(reference Count c);

    private
        type Count = int;

        void init_counter(reference Count c){
            c = 0;
        }

        int get(Count c){
            return c;
        }

        void inc(reference Count c){
            c = c + 1;
        }
}

module Queue{
    public
        type Queue;
        inqueue(reference Queue q, int n);
        int dequeue(reference Queue q);
        ...

    private
        void bookkeep(reference Queue q){
            ...
        }
        ...
}
\end{lstlisting}

\subsection{Moduli Generici}

\begin{lstlisting}
module Buffer<T> imports Counter{
    public
        type Buf;
        void insert(reference Buf f, <T> n);
        <T> get(Buf b);
        Count c;

    private imports Queue<T>{
        type Buf = Queue;

        void insert(reference Buf b, <T> n){
            inqueue(b, n);
            inc(c);
        }

        <T> get(Buf b){
            return dequeue(b);
            inc(c);
        }
    }
}

module Queue<S>{
    public
        type Queue;
        inqueue(reference Queue q, <S> n);
        <S> dequeue(reference Queue q);
        ...

    private
        void bookkeep(reference Queue q){
            ...
        }
        ...
}
\end{lstlisting}

Un modulo generico è uno schema parametrico: il parametro di tipo viene istanziato quando il modulo è usato o collegato.

\section{Limiti degli ADT}

Gli ADT garantiscono \textbf{incapsulamento} e \textbf{information hiding}, ma sono rigidi quando si desidera estendere o riusare un'astrazione.

\subsection{Esempio: Counter}

\begin{lstlisting}
abstype Counter{
    type Counter = int;

    signature
        void reset(Counter x);
        int get(Counter x);
        void inc(Counter x);

    operations
        void reset(Counter x){
            x = 0;
        }

        int get(Counter c){
            return x;
        }

        void inc(Counter c){
            x = x + 1;
        }
}
\end{lstlisting}

Se si vuole definire un contatore arricchito, ad esempio capace di ricordare quante volte è stato azzerato, con gli ADT tradizionali si incontrano due problemi:
\begin{itemize}
  \item duplicare il codice del vecchio ADT nella nuova definizione;
  \item riusare internamente il vecchio ADT, ma senza ereditare automaticamente le sue operazioni.
\end{itemize}

Inoltre, se \texttt{Counter} e \texttt{NewCounter} sono tipi distinti, non è possibile manipolarli uniformemente in una stessa struttura dati senza una nozione di compatibilità.

\subsection{Requisiti che Portano agli Oggetti}

Il paradigma orientato agli oggetti nasce come risposta a quattro esigenze:
\begin{enumerate}
  \item incapsulare dati e operazioni;
  \item ereditare implementazioni già definite;
  \item definire una relazione di compatibilità basata sulle operazioni ammissibili;
  \item selezionare dinamicamente l'operazione corretta in base al tipo effettivo dell'oggetto.
\end{enumerate}

\section{Oggetti e Classi}

\subsection{Oggetti}

Un \textbf{oggetto} è una capsula che contiene:
\begin{itemize}
  \item dati, detti \textbf{instance variables}, \textbf{fields} o \textbf{data members};
  \item operazioni, dette \textbf{methods} o \textbf{member functions};
  \item un'interfaccia visibile all'esterno.
\end{itemize}

L'esecuzione di un metodo avviene tramite invio di un messaggio:

\begin{lstlisting}
object.method(parameters)
\end{lstlisting}

L'oggetto ricevente è un parametro implicito del metodo. Per esempio, \texttt{o.inc()} significa che l'oggetto \texttt{o} riceve il messaggio \texttt{inc}.

\subsection{Classi}

Una \textbf{classe} è un modello per un insieme di oggetti. Stabilisce:
\begin{itemize}
  \item i campi degli oggetti;
  \item la visibilità dei campi;
  \item i metodi disponibili;
  \item la signature, la visibilità e l'implementazione dei metodi.
\end{itemize}

\begin{lstlisting}
class Counter{
    private int x;

    public void reset(){
        x = 0;
    }

    public int get(){
        return x;
    }

    public void inc(){
        x = x + 1;
    }
}
\end{lstlisting}

La creazione dinamica di un oggetto avviene per istanziazione della classe:

\begin{lstlisting}
Counter c = new Counter();
Counter d = new Counter();
\end{lstlisting}

Gli oggetti \texttt{c} e \texttt{d} hanno la stessa struttura, ma campi distinti. Il codice dei metodi può essere memorizzato una sola volta nella classe.

\subsection{Il Riferimento \texttt{this}}

Durante l'esecuzione di un metodo, l'oggetto corrente è accessibile tramite un riferimento implicito, spesso chiamato \texttt{this} o \texttt{self}.

\begin{lstlisting}
public void inc(){
    this.x = this.x + 1;
}
\end{lstlisting}

Le variabili locali sono accedute tramite il record di attivazione; i campi dell'oggetto sono acceduti tramite offset rispetto a \texttt{this}.

\subsection{Classi Statiche e Membri Statici}

Alcuni linguaggi permettono campi e metodi associati alla classe, non alle singole istanze. Essi sono detti \textbf{static} o \textbf{class members}. Un metodo statico non può riferirsi a \texttt{this}, perché non ha un oggetto corrente.

\section{Incapsulamento negli Oggetti}

Ogni classe presenta almeno due viste:
\begin{itemize}
  \item una vista \textbf{private}, completa, disponibile all'interno della classe;
  \item una vista \textbf{public}, limitata, disponibile ai client.
\end{itemize}

La parte pubblica è l'\textbf{interfaccia} della classe. L'incapsulamento degli oggetti è simile a quello degli ADT, ma diventa più flessibile grazie a sottotipi, ereditarietà e dispatch dinamico.

\section{Sottotipi}

In un linguaggio tipato con classi, la definizione di una classe introduce anche un tipo: i suoi valori sono le istanze della classe.

\subsection{Definizione Funzionale di Sottotipo}

Un tipo $T$ è sottotipo di un tipo $S$, scritto:
\[
T <: S
\]
se ogni messaggio compreso dagli oggetti di tipo $S$ è compreso anche dagli oggetti di tipo $T$.

In termini di interfacce:
\[
T <: S \quad \Longleftrightarrow \quad \mathrm{Interface}(S) \subseteq \mathrm{Interface}(T).
\]

Il verso dell'inclusione è importante: un sottotipo ha un'interfaccia più grande o uguale.

\subsection{Subclasses}

Nei linguaggi con equivalenza per nome, la relazione di sottotipo deve spesso essere dichiarata esplicitamente tramite una relazione di \textbf{subclassing}.

\begin{lstlisting}
class NamedCounter extending Counter{
    private String name;

    public void set_name(String n){
        name = n;
    }

    public String get_name(){
        return name;
    }
}
\end{lstlisting}

\texttt{NamedCounter} è sottoclasse di \texttt{Counter}; quindi il tipo \texttt{NamedCounter} è sottotipo di \texttt{Counter}.

\section{Ereditarietà}

L'\textbf{ereditarietà} è il meccanismo che permette a una classe di riusare campi e metodi definiti in una superclasse. Deve essere distinta dalla relazione di sottotipo: spesso coincidono, ma concettualmente sono diverse.

\subsection{Overriding}

Una sottoclasse può ridefinire l'implementazione di un metodo ereditato. Questo meccanismo è detto \textbf{overriding}.

\begin{lstlisting}
class NewCounter extending Counter{
    private int num_reset = 0;

    public void reset(){
        x = 0;
        num_reset = num_reset + 1;
    }

    public int howmany_resets(){
        return num_reset;
    }
}
\end{lstlisting}

Un messaggio \texttt{reset} inviato a un'istanza di \texttt{NewCounter} esegue la nuova implementazione.

\subsection{Shadowing}

Lo \textbf{shadowing} avviene quando una sottoclasse ridefinisce un campo già presente nella superclasse. A differenza dell'overriding dei metodi, lo shadowing è un meccanismo statico.

\begin{lstlisting}
class EvenNewCounter extending NewCounter{
    private int num_reset = 2;

    public void reset(){
        x = 0;
        num_reset = num_reset + 2;
    }

    public int howmany_resets(){
        return num_reset;
    }
}
\end{lstlisting}

I riferimenti a \texttt{num\_reset} dentro \texttt{EvenNewCounter} indicano il campo locale della sottoclasse, non quello della superclasse.

\subsection{Visibilità Protetta}

Molti linguaggi introducono una visibilità \textbf{protected}: un campo o metodo è invisibile ai client generici, ma visibile nelle sottoclassi. Questo rende più efficiente il riuso, ma aumenta l'accoppiamento tra superclasse e sottoclasse.

\subsection{Ereditarietà Singola e Multipla}

Con \textbf{ereditarietà singola}, ogni classe ha al massimo una superclasse immediata; la gerarchia è un albero. Con \textbf{ereditarietà multipla}, una classe può ereditare da più superclassi; la gerarchia diventa un grafo aciclico orientato.

L'ereditarietà multipla introduce conflitti di nomi:

\begin{lstlisting}
class A{
    int x;
    int f(){
        return x;
    }
}

class B{
    int y;
    int f(){
        return y;
    }
}

class C extending A,B{
    int h(){
        return f();
    }
}
\end{lstlisting}

Il riferimento a \texttt{f} in \texttt{C} è ambiguo. Le soluzioni possibili sono:
\begin{enumerate}
  \item proibire sintatticamente i conflitti;
  \item richiedere qualificazione esplicita, ad esempio \texttt{A::f()} o \texttt{B::f()};
  \item adottare una convenzione arbitraria di risoluzione.
\end{enumerate}

\subsection{Diamond Problem}

Il \textbf{diamond problem} nasce quando una classe eredita da due superclassi che condividono una superclasse comune.

\begin{lstlisting}
class Top{
    int w;
    int f(){
        return w;
    }
}

class A extending Top{
    int x;
    int g(){
        return w + x;
    }
}

class B extending Top{
    int y;
    int f(){
        return w + y;
    }
    int k(){
        return y;
    }
}

class Bottom extending A,B{
    int z;
    int h(){
        return z;
    }
}
\end{lstlisting}

Il problema è stabilire se \texttt{Bottom} debba contenere una o due copie della parte ereditata da \texttt{Top}.

\section{Dynamic Method Lookup}

Il \textbf{dynamic method lookup}, o \textbf{dynamic dispatch}, è il cuore del paradigma orientato agli oggetti. Quando un metodo $m$ è invocato su un oggetto $o$:
\[
o.m(\mathrm{parameters})
\]
l'implementazione selezionata dipende dal \textbf{tipo dinamico} dell'oggetto ricevente, non dal tipo statico del riferimento.

\subsection{Esempio di Dispatch Dinamico}

\begin{lstlisting}
class Counter{
    private int x;

    public void reset(){
        x = 0;
    }

    public int get(){
        return x;
    }

    public void inc(){
        x = x + 1;
    }
}

class NewCounter extending Counter{
    private int num_reset = 0;

    public void reset(){
        x = 0;
        num_reset = num_reset + 1;
    }

    public int howmany_resets(){
        return num_reset;
    }
}
\end{lstlisting}

\begin{lstlisting}
NewCounter n = new NewCounter();
Counter c;
c = n;
c.reset();
\end{lstlisting}

Il tipo statico di \texttt{c} è \texttt{Counter}, ma il suo valore dinamico è un'istanza di \texttt{NewCounter}. La chiamata \texttt{c.reset()} invoca quindi il metodo \texttt{reset} ridefinito in \texttt{NewCounter}.

\subsection{Uniformità sui Sottotipi}

\begin{lstlisting}
Counter V[100];

for (int i = 1; i < 100; i = i + 1)
    V[i].reset();
\end{lstlisting}

Ogni elemento dell'array può avere un tipo dinamico diverso, purché sottotipo di \texttt{Counter}. Il dispatch dinamico garantisce che venga eseguito il metodo corretto.

\subsection{Late Binding of Self}

Il riferimento \texttt{this} è legato dinamicamente all'oggetto corrente. Questo produce il cosiddetto \textbf{late binding of self}.

\begin{lstlisting}
class A{
    int a = 0;
    void f(){ g(); }
    void g(){ a = 1; }
}

class B extending A{
    int b = 0;
    void g(){ b = 2; }
}
\end{lstlisting}

Se un'istanza di \texttt{B} invoca \texttt{f}, il corpo ereditato da \texttt{A} chiama implicitamente \texttt{this.g()}; poiché \texttt{this} è l'oggetto di tipo dinamico \texttt{B}, viene invocato \texttt{B.g}.

\subsection{Overriding e Shadowing}

\begin{lstlisting}
class A{
    int a = 1;
    int f(){ return -a; }
}

class B extending A{
    int a = 2;
    int f(){ return a; }
}

B obj_b = new B();
print(obj_b.f());
print(obj_b.a);

A obj_a = obj_b;
print(obj_a.f());
print(obj_a.a);
\end{lstlisting}

Le chiamate di metodo usano dispatch dinamico; l'accesso ai campi usa invece il tipo statico del riferimento. Perciò l'overriding è dinamico, mentre lo shadowing dei campi è statico.

\section{Implementazione degli Oggetti}

\subsection{Rappresentazione degli Oggetti}

Un oggetto può essere rappresentato come un record che contiene:
\begin{itemize}
  \item i campi definiti dalla sua classe;
  \item i campi ereditati dalle superclassi;
  \item eventuali campi omonimi introdotti per shadowing;
  \item un puntatore al descrittore della classe.
\end{itemize}

Se un oggetto di classe $B$ è visto tramite un riferimento statico a una superclasse $A$, il type checking garantisce che siano accessibili solo campi e metodi presenti nell'interfaccia di $A$.

\subsection{Metodo e Parametro \texttt{this}}

Un metodo è eseguito in modo simile a una funzione: viene creato un record di attivazione per parametri, variabili locali e informazioni di controllo. In più, viene passato un puntatore all'oggetto ricevente, cioè \texttt{this}. L'accesso a un campo avviene tramite:
\[
\mathrm{address}(\mathrm{field}) =
\mathrm{this} + \mathrm{offset}(\mathrm{field}).
\]

\subsection{Vtable per Ereditarietà Singola}

In un linguaggio con tipi statici, il dispatch dinamico può essere implementato tramite una \textbf{vtable} (\textit{virtual function table}). Ogni classe ha una tabella contenente i puntatori ai metodi disponibili. La vtable di una sottoclasse:
\begin{itemize}
  \item copia la vtable della superclasse;
  \item sostituisce i metodi ridefiniti;
  \item aggiunge in fondo i nuovi metodi.
\end{itemize}

Se il metodo \texttt{f} è l'$n$-esimo elemento della vtable e ogni indirizzo occupa $w$ byte, una chiamata dinamica può essere schematizzata così:

\begin{lstlisting}
R1 := pa                         // access the object
R2 := *(R1)                      // vtable
R3 := *(R2 + (n - 1) * w)        // indirect to f
call *(R3)                       // call f
\end{lstlisting}

Il costo è costante: due indirezioni più la chiamata.

\subsection{Downcasting}

Il \textbf{downcasting} specializza staticamente un riferimento seguendo la gerarchia dei sottotipi in direzione opposta.

\begin{lstlisting}
Object clone(){...}
\end{lstlisting}

\begin{lstlisting}
A a = (A) o.clone();
\end{lstlisting}

La macchina astratta deve verificare a runtime che il tipo dinamico dell'oggetto sia effettivamente compatibile con \texttt{A}; altrimenti si produce un errore dinamico di tipo.

\subsection{Fragile Base Class Problem}

Il \textbf{fragile base class problem} si verifica quando una modifica apparentemente innocua a una superclasse rompe il comportamento o l'interfaccia binaria delle sottoclassi già compilate.

\begin{lstlisting}
class Upper{
    int x;
    int f(){...}
}

class Lower extending Upper{
    int y;
    int g(){ return y + f(); }
}
\end{lstlisting}

Se la superclasse viene modificata:

\begin{lstlisting}
class Upper{
    int x;
    int h(){ return x; }
    int f(){...}
}
\end{lstlisting}

gli offset dei metodi nella vtable possono cambiare. Una sottoclasse compilata in precedenza può quindi accedere al metodo sbagliato, a meno di ricompilazione o risoluzione dinamica degli offset.

\subsection{Dispatch nella JVM}

La Java Virtual Machine usa un meccanismo basato su \textbf{constant pool}, caricamento dinamico delle classi e risoluzione lazy dei nomi. Una chiamata virtuale può inizialmente essere rappresentata come:

\begin{lstlisting}
invokevirtual index
\end{lstlisting}

Alla prima esecuzione, la JVM risolve il nome, controlla visibilità e tipi, calcola l'offset nella tabella dei metodi e può riscrivere l'istruzione in forma ottimizzata:

\begin{lstlisting}
invokevirtual_quick d np
\end{lstlisting}

Per i metodi di interfaccia, l'offset può variare tra classi diverse che implementano la stessa interfaccia:

\begin{lstlisting}
interface Interface{
    void foo();
}

public class A implements Interface{
    int x;
    void foo(){...}
}

public class B implements Interface{
    int y;
    int fie(){...}
    int foo(){...}
}
\end{lstlisting}

La JVM deve mantenere anche informazioni sul nome e sulla signature del metodo:

\begin{lstlisting}
invokeinterface_quick name_of_foo, d
\end{lstlisting}

\section{Implementazione dell'Ereditarietà Multipla}

\subsection{Struttura delle Vtable}

\begin{lstlisting}
class A{
    int x;
    int f(){
        return x;
    }
}

class B{
    int y;
    int g(){
        return y;
    }
    int h(){
        return 2 * y;
    }
}

class C extending A,B{
    int z;
    int g(){
        return x + y + z;
    }
    int k(){
        return z;
    }
}
\end{lstlisting}

In presenza di ereditarietà multipla, un oggetto può avere più viste statiche. Una istanza di \texttt{C} deve poter essere vista come \texttt{A}, come \texttt{B} e come \texttt{C}. La rappresentazione contiene più sezioni e, spesso, più puntatori a vtable.

Una chiamata a un metodo ereditato dalla seconda superclasse richiede una correzione del puntatore \texttt{this}. Se \texttt{h} è l'$n$-esimo metodo nella vtable di \texttt{B}, uno schema possibile è:

\begin{lstlisting}
R1 := pa                              // view A
R1 := R1 + d                          // view B
R2 := *(R1)                           // vtable for B
R3 := *(R2 + (n - 1) * 2 * w)         // address of h
R2 := *(R2 + (n - 1) * 2 * w + w)     // correction
this := R1 + R2
call *(R3)                            // call h
\end{lstlisting}

\subsection{Ereditarietà Multipla Replicata}

Nell'ereditarietà multipla \textbf{replicata}, una classe in fondo al diamante contiene copie distinte della superclasse comune. Questo semplifica l'implementazione, ma rende ambigui accessi e conversioni verso la superclasse condivisa.

\subsection{Ereditarietà Multipla Condivisa}

Nell'ereditarietà multipla \textbf{condivisa}, la classe in fondo al diamante contiene una sola copia della superclasse comune. In C\texttt{++}, questo comportamento è ottenuto tramite \textbf{virtual base classes}.

L'accesso alla parte condivisa richiede un'indirezione aggiuntiva:

\begin{lstlisting}
R1 := pa                              // view Bottom
R1 := *(R1 + w)                       // view Top
R2 := *(R1)                           // vtable for Top
R3 := *(R2 + (n - 1) * 2 * w)         // address of f
R2 := *(R2 + (n - 1) * 2 * w + w)     // correction
this := R1 + R2
call *(R3)
\end{lstlisting}

\section{Polimorfismo per Sottotipo}

Il \textbf{polimorfismo per sottotipo} è una forma di polimorfismo universale limitata ai sottotipi di un tipo dato. Se un metodo ha forma:

\begin{lstlisting}
B foo(A x){...}
\end{lstlisting}

allora può ricevere un valore di qualunque sottotipo di \texttt{A}. In notazione formale:
\[
\forall T <: A.\ T \to B.
\]

\subsection{Identità e Limiti del Type System}

\begin{lstlisting}
A Ide(A x){
    return x;
}
\end{lstlisting}

Semanticamente, il tipo più preciso sarebbe:
\[
\forall T <: A.\ T \to T.
\]

Tuttavia molti linguaggi non sono in grado di esprimere staticamente che il tipo del risultato dipende dal tipo dell'argomento. Perciò può essere necessario un cast:

\begin{lstlisting}
C cc = (C) Ide(c);
\end{lstlisting}

\subsection{Stack Basato su \texttt{Object}}

\begin{lstlisting}
class Elem{
    Object info;
    Elem prox;
}

class Stack{
    private Elem top = null; // empty stack

    boolean isEmpty(){
        return top == null;
    }

    void push(Object o){
        Elem ne = new Elem();
        ne.info = o;
        ne.prox = top;
        top = ne;
    }

    Object pop(){
        Object tmp = top.info;
        top = top.prox;
        return tmp;
    }
}
\end{lstlisting}

Poiché \texttt{pop} restituisce \texttt{Object}, l'uso specifico richiede un cast dinamico:

\begin{lstlisting}
C c;
Stack s = new Stack();
s.push(new C());
c = (C) s.pop(); // dynamic check
\end{lstlisting}

\section{Generics in Java}

I \textbf{generics} introducono polimorfismo parametrico esplicito in Java. Classi, interfacce e metodi possono avere parametri di tipo.

\subsection{Stack Generico}

\begin{lstlisting}
class Elem<A>{
    A info;
    Elem<A> prox;
}

class Stack<A>{
    private Elem<A> top = null; // empty stack

    boolean isEmpty(){
        return top == null;
    }

    void push(A o){
        Elem<A> ne = new Elem<A>();
        ne.info = o;
        ne.prox = top;
        top = ne;
    }

    A pop(){
        A tmp = top.info;
        top = top.prox;
        return tmp;
    }
}
\end{lstlisting}

L'uso diventa type-safe senza cast espliciti:

\begin{lstlisting}
Stack<String> ss = new Stack<String>();
ss.push(new String("pippo"));
String s = ss.pop();
\end{lstlisting}

\subsection{Coppie Generiche}

\begin{lstlisting}
class Pair<A,B>{
    private A a;
    private B b;

    Pair(A x, B y){
        a = x;
        b = y;
    }

    A First(){
        return a;
    }

    B Second(){
        return b;
    }
}
\end{lstlisting}

\begin{lstlisting}
Integer i = new Integer(3);
String v = new String("pippo");
Pair<Integer,String> c = new Pair<Integer,String>(3, v);
String w = c.Second();
\end{lstlisting}

\subsection{Metodi Generici}

\begin{lstlisting}
Pair<Object,Object> diagonal(Object x){
    return new Pair<Object,Object>(x, x);
}
\end{lstlisting}

Questa versione perde informazione sul tipo specifico dell'argomento. Una versione generica mantiene la relazione tra argomento e risultato:

\begin{lstlisting}
<T> Pair<T,T> diagonal(T x){
    return new Pair<T,T>(x, x);
}
\end{lstlisting}

Formalmente:
\[
\forall T.\ T \to Pair<T,T>.
\]

Esempi di inferenza del parametro:

\begin{lstlisting}
Pair<Integer,Integer> ci = diagonal(new Integer(4));
Pair<String,String> cs = diagonal(new String("pippo"));
\end{lstlisting}

\subsection{Bounded Generics}

Un parametro di tipo può essere vincolato a un sottotipo:

\begin{lstlisting}
interface Shape{
    void draw();
}

class Circle implements Shape{
    public void draw(){...}
}

class Rhombus implements Shape{
    public void draw(){...}
}
\end{lstlisting}

\begin{lstlisting}
<T extends Shape> void drawAll(List<T> forms){
    for(Shape f : forms)
        f.draw();
}
\end{lstlisting}

Il tipo formale è:
\[
\forall T <: Shape.\ List<T> \to void.
\]

Il metodo accetta liste di elementi appartenenti a qualunque sottotipo di \texttt{Shape}.

\subsection{F-Bounded Polymorphism}

Un parametro di tipo può comparire nel proprio vincolo. Questo schema è detto \textbf{F-bounded polymorphism}.

\begin{lstlisting}
public static <T extends Comparable<T>>
T max(List<T> list)
\end{lstlisting}

Questa firma richiede che $T$ sia confrontabile con valori dello stesso tipo. Una forma più flessibile consente il confronto con supertipi di $T$:

\begin{lstlisting}
public static <T extends Comparable<? super T>>
T max(List<T> list)
\end{lstlisting}

\subsection{Wildcard}

Il carattere \texttt{?} indica un tipo sconosciuto. \texttt{List<?>} è un supertipo comune di tutte le liste \texttt{List<A>}.

\begin{lstlisting}
void drawAll(List<? extends Shape> forms){
    for(Shape f : forms)
        f.draw();
}
\end{lstlisting}

Una wildcard è utile quando il parametro di tipo compare una sola volta. Se compare più volte, un parametro esplicito rende il vincolo più chiaro.

\subsection{Implementazione dei Generics in Java}

Java implementa i generics tramite \textbf{erasure}. Dopo il controllo statico:
\begin{itemize}
  \item le informazioni tra parentesi angolari sono rimosse;
  \item le variabili di tipo sono sostituite dal loro upper bound, spesso \texttt{Object};
  \item il compilatore inserisce cast dinamici dove necessario.
\end{itemize}

Conseguenze:
\begin{itemize}
  \item la JVM non deve essere modificata;
  \item codice generico e non generico possono coesistere;
  \item eventuali buchi locali nel sistema statico sono intercettati a runtime dai cast inseriti.
\end{itemize}

\section{Generics, Array e Gerarchia dei Sottotipi}

In Java, i generics non preservano la gerarchia dei sottotipi:
\[
A <: B \quad \nRightarrow \quad \mathrm{Stack}\langle A\rangle <: \mathrm{Stack}\langle B\rangle.
\]

\subsection{Perché i Generics non sono Covarianti}

\begin{lstlisting}
Stack<Integer> si = new Stack<Integer>();
Stack<Object> so = si; // Incorrect in Java
\end{lstlisting}

Se la seconda riga fosse ammessa, si potrebbe scrivere:

\begin{lstlisting}
so.push(new String("pluto"));
Integer i = si.pop(); // danger!
\end{lstlisting}

Il risultato sarebbe una violazione della sicurezza dei tipi: una stringa verrebbe estratta da una pila che staticamente promette interi.

\subsection{Array Covarianti}

Gli array Java sono invece covarianti:
\[
A <: B \Rightarrow A[] <: B[].
\]

\begin{lstlisting}
Integer[] ai = new Integer[10];
Object[] ao = ai;              // correct: Integer[] <: Object[]
ao[0] = new String("pluto");   // correct statically; runtime error
\end{lstlisting}

La sicurezza è mantenuta tramite controlli dinamici sugli aggiornamenti dell'array.

\subsection{Covarianza e Contravarianza}

Sia $(D,\leq)$ un insieme preordinato. Una funzione $f : D \to D$ è:
\begin{itemize}
  \item \textbf{covariante} se $x \leq y \Rightarrow f(x) \leq f(y)$;
  \item \textbf{contravariante} se $x \leq y \Rightarrow f(x) \geq f(y)$.
\end{itemize}

Nel contesto dei tipi, la terminologia descrive come un costruttore di tipo preserva o inverte la relazione di sottotipo.

\section{Overriding Covariante e Contravariante}

\subsection{Tipo del Risultato}

Supponiamo una classe:

\begin{lstlisting}
class F{
    C fie(A p){...}
}
\end{lstlisting}

Se $D <: C$, una sottoclasse può ridefinire il metodo con risultato più specifico:

\begin{lstlisting}
class E extending F{
    D fie(A p){...}
}
\end{lstlisting}

La ridefinizione è corretta perché ogni valore di tipo $D$ è anche un valore di tipo $C$. Quindi l'overriding è \textbf{covariante} nel tipo del risultato.

\subsection{Tipo degli Argomenti}

Per gli argomenti, la regola semanticamente corretta è opposta. Un metodo:
\[
S \to T
\]
è sottotipo di:
\[
S' \to T'
\]
quando:
\[
S' <: S
\qquad \text{e} \qquad
T <: T'.
\]

La funzione è quindi \textbf{contravariante} negli argomenti e \textbf{covariante} nel risultato.

\subsection{Binary Methods}

La contravarianza degli argomenti è spesso controintuitiva nei \textbf{binary methods}, cioè metodi che ricevono un parametro dello stesso tipo dell'oggetto ricevente.

\begin{lstlisting}
class Point{
    boolean eq(Point p){...}
}

class ColoredPoint extending Point{
    boolean eq(ColoredPoint p){...}
}
\end{lstlisting}

Secondo la regola contravariante, \texttt{ColoredPoint.eq} non è una ridefinizione type-safe di \texttt{Point.eq}, perché restringe il tipo dell'argomento invece di allargarlo. Per questo molti linguaggi, come Java e C\texttt{++}, richiedono identità dei tipi degli argomenti nei metodi ridefiniti.

\section{Linguaggi Basati su Delegation}

Non tutti i linguaggi orientati agli oggetti sono basati su classi. Nei linguaggi a \textbf{delegation} o \textbf{prototype-based}, un oggetto può delegare a un altro oggetto, detto \textbf{parent}, l'esecuzione di un metodo che non possiede.

Un oggetto può essere creato clonando un \textbf{prototype}. Se riceve un messaggio che non comprende, lo inoltra al parent. Il riferimento a \texttt{self} resta legato all'oggetto originario che ha ricevuto il messaggio.

Questo modello rende più flessibile l'ereditarietà: un oggetto può cambiare dinamicamente il proprio parent e quindi il proprio comportamento.

\section{Sintesi Concettuale}

\subsection{Astrazione dei Dati}

L'astrazione dei dati si fonda su:
\begin{itemize}
  \item separazione tra interfaccia e implementazione;
  \item opacità della rappresentazione;
  \item specificazione del comportamento osservabile;
  \item indipendenza dalla rappresentazione;
  \item modularizzazione per sistemi grandi.
\end{itemize}

\subsection{Paradigma Orientato agli Oggetti}

Il paradigma orientato agli oggetti combina quattro meccanismi:
\begin{itemize}
  \item \textbf{incapsulamento} dei dati;
  \item \textbf{sottotipi} come relazione di compatibilità;
  \item \textbf{ereditarietà} come riuso di implementazione;
  \item \textbf{dynamic dispatch} come selezione runtime dei metodi.
\end{itemize}

\subsection{Confronto Finale}

\begin{center}
\begin{tabular}{|p{3.3cm}|p{5.2cm}|p{5.2cm}|}
\hline
\textbf{Meccanismo} & \textbf{Idea centrale} & \textbf{Limite o costo} \\
\hline
ADT & Incapsula un tipo e le sue operazioni & Poco flessibile nell'estensione \\
\hline
Modulo & Incapsula più dichiarazioni e controlla la visibilità & Dipende spesso da meccanismi di compilazione separata \\
\hline
Classe & Modello per oggetti con campi e metodi & Introduce complessità di dispatch e rappresentazione \\
\hline
Sottotipo & Compatibilità basata sull'interfaccia & Può divergere dall'ereditarietà \\
\hline
Ereditarietà & Riuso di implementazione & Fragile base class e conflitti multipli \\
\hline
Dynamic dispatch & Selezione runtime del metodo & Richiede supporto nella macchina astratta \\
\hline
Generics & Polimorfismo parametrico esplicito & Interazione delicata con sottotipi e array \\
\hline
\end{tabular}
\end{center}

La traiettoria concettuale è chiara: dagli ADT, che proteggono una rappresentazione dietro una signature, si arriva agli oggetti, che rendono l'astrazione estensibile tramite sottotipi, ereditarietà e dispatch dinamico.

\end{document}
